\section{Conclusion}
\label{sec:conclusion}

Reliable, real-time traffic anomaly detection in intelligent transportation systems is frequently hindered by decentralized data silos, extreme class imbalance, and the strict hardware constraints of vehicular edge nodes. To address these challenges, this paper presented FedRouteAI, a federated learning framework integrating a deployment-aware inference pipeline with an adaptive aggregation strategy. The core methodological contributions include a federated architecture that preserves data privacy across distinct regional clients, a historical feature engineering module for localized normalization, the four-factor AdaptiveFedAvg aggregation algorithm, and a deployment-aware inference mechanism relying on zero-imputation and dynamic Out-of-Distribution (OOD) warnings.

The experimental evaluation indicates that, under the evaluated synthetic dataset, the federated approaches achieved higher generalization performance than the centralized baseline. Furthermore, the results suggest that AdaptiveFedAvg achieves performance comparable to standard FedAvg, with only marginal differences in global accuracy and Macro F1 scores, indicating that standard, dataset-proportional aggregation remains robust in stable deployment environments. Additionally, systems-level benchmarking shows that the lightweight Multi-Layer Perceptron architecture operates with low inference latency and minimal communication overhead, indicating that the proposed framework satisfies the latency and communication requirements evaluated in this study and is compatible with real-time edge deployment scenarios.

Despite these promising indications, several limitations must be acknowledged. The experimental evaluation relied entirely on a synthetic dataset and a simulated federated environment, which limits the validation of the system under true physical conditions. Furthermore, the evaluation was constrained to a single MLP architecture and did not test the anomaly warning layer against the severe, erratic packet drops typical of physical edge sensors. Future work should also evaluate the framework using larger and more heterogeneous federated deployments, investigate more advanced learning architectures (e.g., Graph Neural Networks or Transformer-based models), and validate the system under real roadside infrastructure.
